% =====================================================================
% Rubric Deliverable P2 — Non-AI Baseline Rigor (with numbers)
% MRI-ReportGenerator (EECE503N). Self-contained; compile with pdfLaTeX.
% Every figure below is a primary-source value; authors were verified against PMIDs.
% =====================================================================
\documentclass[11pt]{article}
\usepackage[margin=1in]{geometry}
\usepackage{booktabs}
\usepackage{tabularx}
\usepackage{array}
\usepackage{enumitem}
\usepackage{xcolor}
\usepackage{textcomp}
\usepackage[hidelinks]{hyperref}

\newcolumntype{Y}{>{\raggedright\arraybackslash}X}
\setlist{nosep, leftmargin=1.4em}
\sloppy

\title{\textbf{The Non-AI Baseline: Slow and Unreliable}\\[2pt]
\large Rubric Deliverable P2 --- MRI-ReportGenerator (Cervical Spine MRI Analysis)}
\author{EECE503N Final Project}
\date{2026-06-08}

\begin{document}
\maketitle

\noindent\textbf{Claim.} The non-AI baseline our system augments is \emph{manual radiologist
measurement} of cervical spine MRI --- a radiologist visually inspects the sagittal T2 series and
hand-measures vertebral, disc, canal, and cord geometry. It is the current standard of care, and it is
both \textbf{slow} and \textbf{unreliable}, with inter-observer agreement collapsing for precisely the
structures we automate. Below we quantify both costs from primary sources (PMIDs given; authors verified
against their PMID). Citations carry honest provenance flags: \emph{[cervical MRI]} where directly on
point, and \emph{[borrow]} where the only available figure is from radiograph/CT or the lumbar spine.

% ---------------------------------------------------------------------
\section{Cost 1 --- Time (slow)}
\begin{itemize}
  \item \textbf{Reading time.} Cervical MRI reporting takes \(\approx\) \textbf{2.7\,min median /
  3.8\,min mean} of PACS interaction (excluding dictation) --- Forsberg et al., 2017 (PMID~27714473).
  \item \textbf{Measurement time.} \textbf{No formal time-motion study} isolates the hand-geometry
  measurement time for cervical MRI (a documented negative); the only figure is an informal
  \textbf{5--10\,min/case} self-report --- Zhu et al., 2024 (PMID~38269650, discussion). The closest
  \emph{formal} comparator is full-spine radiograph alignment measurement at
  \(147\pm19\)\,s (expert) rising to \(232\pm20\)\,s (resident) --- Pucciarelli et al., 2025
  (PMID~41003360) \emph{[borrow]}.
\end{itemize}
\noindent A deterministic pipeline computes the same geometry in \(\approx\)\,seconds per case (measured
per-component latency; see the deployment/tradeoffs docs), removing the per-case minutes spent on
repetitive caliper placement.

% ---------------------------------------------------------------------
\section{Cost 2 --- Inter-observer variability (unreliable)}
The decisive cost is reproducibility. Where two radiologists (or the same radiologist on different days)
measure the same study, they disagree --- and the disagreement is \emph{worst} for the cord, compression,
and disc-degeneration readings that drive myelopathy/radiculopathy decisions. A deterministic pipeline has
inter-observer ICC \(=1.0\) by construction.

\begin{table}[h]
\centering\footnotesize
\begin{tabularx}{\textwidth}{@{}>{\hsize=1.25\hsize}Y >{\hsize=0.95\hsize}Y >{\hsize=0.80\hsize}Y@{}}
\toprule
\textbf{Measurement (cervical MRI unless flagged)} & \textbf{Human inter-observer agreement} & \textbf{Source (PMID)} \\
\midrule
\textbf{Cord AP diameter} & ICC \textbf{0.66} (axial) / 0.82 (mid-sag) & Grochmal 2018 (29913296) \\
\textbf{Cord compression (trauma)} & ICC \textbf{0.35--0.56} & Fehlings 2006 (16816769) \\
Cord compression ratio / MSCC & ICC 0.79--0.86 & Karpova 2013 (22772577); 2010 (23637669) \\
\textbf{Disc-degeneration grade ($\kappa$)} & \textbf{0.265} (Pfirrmann-on-cervical); 0.364 (Miyazaki, replication) & Urbanschitz 2021 (34966859); Otaki 2021 (34803082) \\
\textbf{Cobb / alignment} & inter 0.88 / intra 0.83; \textbf{mixed-reader \(\sim\)0.55} & Sevin 2025 (41011045); Lee 2017 (28574884) \\
Cobb (radiograph borrow) & inter 0.886 / intra 0.958, SEM 1.8\textdegree, MDC 5.0\textdegree & Marques 2019 (31668048) \emph{[borrow]} \\
Canal AP diameter & ICC 0.72--0.99 & R\"uegg 2015 (25525959); Griffith 2024 (38617159) \\
Canal-stenosis grade & ICC 0.72--0.80, $\kappa$ 0.60--0.62 & Kang 2011 (21700974) \\
Disc height (mid-disc) & ICC 0.91--0.95 & van Santbrink 2026 (41567758) \\
Vertebral-body height & \emph{no cervical-MRI ICC exists}; lumbar 0.97 & Germann 2022 (36576545) \emph{[borrow]} \\
Space available for cord (SAC, mm) & \emph{no primary cervical ICC exists} & --- (negative) \\
\bottomrule
\end{tabularx}
\end{table}

\noindent\textbf{Lead with where agreement collapses:} cord AP axial ICC \textbf{0.66}, cord compression
ICC \textbf{0.35--0.56}, Pfirrmann-on-cervical $\kappa$ \textbf{0.265}, and Cobb dropping from
\(\sim\)0.96 (expert pairs) to \(\sim\)\textbf{0.55} across mixed readers. These are not edge cases ---
they are the measurements that decide whether a cord is compressed and a disc degenerated.

% ---------------------------------------------------------------------
\section{Why the baseline is insufficient, and why our approach is justified}
\begin{itemize}
  \item \textbf{It does not scale.} Minutes of specialist time per case on repetitive geometry, at volume.
  \item \textbf{It is not reproducible.} The same study yields materially different numbers between
  readers (Section~2); a threshold-based decision sitting on a measurement with ICC 0.35--0.66 inherits
  that noise.
  \item \textbf{Our (non-LLM) AI approach removes both costs by construction.} Automated segmentation
  (TotalSpineSeg / Spinal Cord Toolbox / SPINEPS) + deterministic geometric measurement gives
  \emph{perfectly reproducible} numbers in seconds, and every flag is tied to a cited threshold (no
  black-box verdict). It is consistency + speed + auditability, not a diagnosis --- outputs are
  ``flagged for physician review.''
\end{itemize}

% ---------------------------------------------------------------------
\section{Honest provenance (the negatives are themselves the argument)}
Several baseline figures are \emph{borrows} (vertebral-body-height ICC is lumbar; some Cobb/disc figures
are radiograph/CT), flagged above. More tellingly, for \textbf{vertebral-body height} and \textbf{SAC in
millimetres} \emph{no cervical-MRI inter-observer study exists at all}. These gaps cut in our favour: the
manual baseline is so under-characterised that, for some of the exact measurements our pipeline produces
deterministically, the literature cannot even state how much two humans would disagree. We automate what
has not been reliably measured by hand.

% ---------------------------------------------------------------------
\begin{thebibliography}{9}
\bibitem{forsberg} Forsberg D, et al. Cervical spine MRI reporting time. 2017. PMID:~27714473.
\bibitem{zhu} Zhu, et al. Self-reported manual cervical-MRI measurement effort. 2024. PMID:~38269650.
\bibitem{pucciarelli} Pucciarelli, et al. Spinal alignment measurement time (radiograph). 2025. PMID:~41003360.
\bibitem{grochmal} Grochmal, et al. Cord AP diameter reliability (cervical MRI). 2018. PMID:~29913296.
\bibitem{fehlings} Fehlings MG, et al. Cord-compression measurement reliability. 2006. PMID:~16816769.
\bibitem{karpova} Karpova, et al. Compression-ratio / MSCC reliability. 2013/2010. PMID:~22772577 / 23637669.
\bibitem{urbanschitz} Urbanschitz, et al. Pfirrmann-on-cervical grading $\kappa$. 2021. PMID:~34966859.
\bibitem{otaki} Otaki, et al. Miyazaki grading replication $\kappa$. 2021. PMID:~34803082.
\bibitem{sevin} Sevin, et al. Cervical-MRI Cobb reliability. 2025. PMID:~41011045.
\bibitem{lee2017} Lee, et al. Mixed-reader Cobb spread. 2017. PMID:~28574884.
\bibitem{marques} Marques, et al. Cobb reliability (radiograph; corrected values). 2019. PMID:~31668048.
\bibitem{ruegg} R\"uegg, et al. Canal AP reliability. 2015. PMID:~25525959.
\bibitem{kang} Kang, et al. Canal-stenosis grading reliability. 2011. PMID:~21700974.
\bibitem{santbrink} van Santbrink, et al. Cervical disc-height reliability (MRI). 2026. PMID:~41567758.
\bibitem{germann} Germann, et al. Vertebral-body-height reliability (lumbar MRI borrow). 2022. PMID:~36576545.
\end{thebibliography}

\end{document}
